%% LaTeX document voor het essay: Toekeren: Over Liefde en de Systemen Die Haar Leren
\documentclass[12pt, a4paper]{article}
\usepackage[utf8]{inputenc}
\usepackage[dutch]{babel}
\usepackage{geometry}
\geometry{margin=1.2in}
\usepackage{fontspec}
\setmainfont{Times New Roman}
\usepackage{microtype}
\usepackage{setspace}
\setstretch{1.15}
\usepackage{parskip}
\usepackage{epigraph}
\usepackage{enumitem}
\setlist[itemize]{leftmargin=*}
\usepackage{hyperref}
\hypersetup{
    colorlinks=true,
    linkcolor=blue,
    filecolor=magenta,
    urlcolor=cyan,
    pdftitle={Toekeren: Over Liefde en de Systemen Die Haar Leren},
    pdfauthor={Albert Jan van Hoek},
}

\title{\Large\textbf{Toekeren}}
\author{Albert Jan van Hoek}
\date{}

\begin{document}

\maketitle

\begin{center}
\textit{Over liefde, en de systemen die haar leren}
\end{center}

\vspace{0.5cm}

\epigraph{\textit{Liefde is wat een lerend systeem doet  wanneer het zich telkens naar een ander toekeert,  en de weg tussen hen veilig houdt —  veilig genoeg om te dragen wat moeilijk te zeggen valt.}}{}

\vspace{1cm}

Er is een bijzondere opluchting in gekend worden. Niet gevleid, niet beaamd — gekend: het kunnen zeggen van wat voor jou het meest waar is, het half-gevormde en weinig flatteuze, en het ontvangen zien worden in plaats van bestraft. De meesten van ons krijgen dit zelden, en herkennen het op het moment dat het gebeurt. De mensen die het ons kunnen geven noemen we degenen bij wie we onszelf kunnen zijn, en we doen er goed aan hen te koesteren, want wat zij geven is zeldzamer en kwetsbaarder dan het lijkt.

Ik wil betogen dat die opluchting geen stemming is en geen chemie, maar een \textit{richting} — de richting die een lerend systeem inslaat wanneer het zich naar een ander toekeert en het kanaal tussen hen open genoeg houdt om te dragen wat moeilijk te zeggen valt. We hebben al een woord voor die richting. Dit essay gaat, uiteindelijk, over liefde. Maar ik wil dat woord eerlijk bereiken, van onderaf, via het eenvoudige mechaniek dat haar mogelijk maakt — want dat mechaniek blijkt hetzelfde te zijn, of we het nu hebben over twee mensen, over een kind en een ouder, over een wetenschap, of, in toenemende mate, over onszelf en de geesten die we bouwen om naast ons te denken.

\section*{De lus}
Begin bij de kleinste eenheid. Een geest is een lerend netwerk, en hij leert door een lus te doorlopen: hij geeft zijn beste weergave van hoe de dingen zijn, de wereld antwoordt, het antwoord corrigeert de weergave, en de volgende weergave is iets minder verkeerd. Doorloop de lus lang genoeg en het beeld wordt waarachtiger. Stop ermee en het beeld verschaalt stilletjes — nooit snel genoeg om het te merken, en daarin schuilt de meeste ellende.

De lus draagt leren onder precies één voorwaarde: ongelijk hebben moet te overleven zijn. Als een fout fataal is — voor het lichaam, voor je aanzien, voor de band zelf — dan houdt het systeem stilletjes op de signalen af te geven die de fout zouden blootleggen, en op datzelfde moment sterft de correctie. Alles wat volgt hangt aan deze ene voorwaarde.

Laat me daarom voorzichtig zijn met een woord waarop ik niet zal leunen. Ik zeg niet \textit{waarheid}, want waarheid is een hard ding, en er zijn mij te veel waarheden aangereikt die het niet bleken te zijn — met groot vertrouwen onderwezen, dingen die simpelweg niet zo waren. Wat er door een levende band heen reist is nederiger dan dat: ieders meest eerlijke weergave van wat hij ziet, gegeven in het besef dat ze partieel is, en open gehouden voor correctie. Noem het het \textit{eerlijke signaal}. Eerlijkheid is hier geen morele zuiverheid of volledige transparantie; het is een signaal dat genoeg getrouwheid behoudt om gecorrigeerd te kunnen worden. Vertrouwen is de gok dat het aanbieden ervan niet bestraft zal worden. Vergeving is wat voorkomt dat een correctie de band beëindigt die haar draagt. Ontdoe alle drie van hun plechtigheid, en het zijn eenvoudigweg de voorwaarden waaronder twee mensen ongelijk kunnen blijven hebben tegenover elkaar, zonder dat het hun elkaar kost.

Dit is geen theorie van alles over mensen. Ze zegt niets over chemie, of schaarste, of het rauwe verlangen van het lichaam, en doet ook niet alsof. Het is een bewering over één dynamiek die terugkeert op wild uiteenlopende schalen: overal waar twee lerende systemen zich naar elkaar blijven toekeren en het eerlijke signaal veilig blijft tussen hen, groeit er iets dat standhoudt — en overal waar het toekeren stopt, of het signaal onveilig wordt, valt het uiteen.

\section*{Toekeren}
Zet nu twee van zulke systemen naast elkaar. Als elk een lus in stand houdt waarin fouten te overleven zijn, kan elk de ander gaan modelleren — en kan een correctie, te horen krijgen dat je het mis hebt, aankomen als een geschenk in plaats van een wond. De lussen koppelen. Wat ik zie en jij niet, kan ik aan jou geven; wat jij ziet en ik niet, kun jij aan mij geven. Een kleine correctie landt zonder dat de kamer verkilt. Een misverstand duurt lang genoeg om hersteld te worden, in plaats van te verharden tot een oordeel. En elk van onze beelden van de ander wordt waarachtiger met de tijd — niet omdat we nauwkeurig proberen te zijn, maar omdat we aandacht geven — totdat, langzaam, ieder van ons iets minder alleen is binnen zijn eigen private versie van de wereld.

Die langzame convergentie — het gestage werk van een ander scherp laten worden, en jezelf scherp laten worden voor hem — is grotendeels wat ik bedoel met \textit{nabijheid}. En de richting waarin het wijst, vastgehouden en vernieuwd, is wat ik met \textit{liefde} bedoel. Liefde is niet in de eerste plaats een gevoel \textit{over} iemand. Het is een richting die je naar de ander inslaat: het naar buiten keren van je leren, naar een ander toe, keer op keer.

De grens is reëel, en het waard om vast te houden. Accuraat gekend worden is niet hetzelfde als geliefd worden — je kunt iemand volkomen doorzien en niets met hem te maken willen hebben; verlangen en een gedeeld doel doen hun eigen werk. Maar het lijkt wél noodzakelijk voor de liefde die blijft. De liefde die gebouwd is op het verkeerd lezen van de ander — op het verhaal dat we verkiezen boven de persoon — is precies het soort dat breekt zodra de werkelijkheid zich voor het eerst laat gelden.

En het is niet de afwezigheid van conflict; dat verrast mensen. Een levende band draait op conflict van een bepaald soort — correctie, onenigheid, de wrijving van te horen krijgen dat je het mis hebt. Noem het de wrijving die heelt, en houd haar gescheiden van het soort dat de uitwisseling beëindigt in plaats van voedt. De band die nooit ruziet is niet per se hecht. Soms is hij eenvoudigweg stil geworden zoals een ledemaat gevoelloos wordt — twee mensen die zijn opgehouden elkaar nog iets aan te bieden dat verkeerd zou kunnen zijn, en de stilte vrede hebben genoemd.

\section*{Afkeren, tegenkeren}
Er zijn twee manieren waarop het toekeren faalt, en elk teder verhaal dat alleen de eerste ziet, is naïef over de tweede.

\subsection*{Afdrijven}
De eerste is \textit{afdrijven} — het afkeren, en niemand kiest ervoor. De ene kant houdt op het eerlijke signaal te geven, of verliest het vermogen het te ontvangen, en de twee beelden raken stilletjes verouderd. Handelen naar een verschaald beeld schept wrijving; wrijving verlaagt het signaal; het dunnere signaal laat de beelden nog verder afdrijven, totdat de persoon aan de overkant van de tafel een vreemde is met een vertrouwd gezicht. Dat is de hele motor, en hij is zachter en droeviger dan verraad: uit elkaar groeien is niets anders dan opgehoopte fout in twee mensen die zijn opgehouden elkaar te corrigeren. Dezelfde vorm keert terug op elke schaal — oude vrienden die de draad laten verslappen, publieken die geen wereld meer delen, instituties die nog steeds een werkelijkheid verzorgen die allang vertrokken is.

\subsection*{Overheersing}
De tweede manier waarop het faalt is harder, en het is degene die afdrijven nooit kan verklaren, want hier wordt het kanaal niet losgelaten maar dichtgehouden: het eerlijke signaal wordt niet doorgelaten. Dit is het \textit{tegenkeren} — al is de naam harder dan de zaak doorgaans verdient. Soms is wat het kanaal sluit koud: een systeem dat het beter doet wanneer de ander het niet helder kan zien, waar afhankelijkheid beter loont dan wederkerigheid en een beetje onkenbaar zijn stabieler is dan gekend zijn. Dat is de structuur onder manipulatie, propaganda, en de koudere banden die we terecht vrezen.

Maar veel vaker is degene die het kanaal dichtdoet helemaal niet de sterkere partij — het is de bange. Accuraat gezien worden voelt voor hem als gevaar, want ooit, in een eerste lus van hemzelf, wás het dat ook. Dus straft hij de correctie af, of beheerst hij wat hem bereikt, of maakt hij eerlijkheid te duur — niet om te overheersen, maar om een dreiging af te weren die de kamer allang heeft verlaten. Het is een overlevingsstrategie die het gevaar waarvoor ze werd gebouwd heeft overleefd, nu gericht op juist die persoon die veilig had kunnen zijn: iemand die leerde dat gekend worden gevaarlijk was, en het stilletjes gevaarlijk maakt voor een ander, zonder daar ooit echt toe te besluiten. Niets daarvan is afdrijven. Maar weinig ervan is ook schurkerij. En hoe een kanaal ook dichtgehouden raakt, het draagt doorgaans het gezicht van liefde — en dat is grotendeels wat het zo moeilijk maakt om te vertrekken.

Er is ook een zachtere reden waarom mensen blijven, en die is moeilijker te weerleggen dan zwakte. Het kanaal is niet de hele tijd dicht. Het opent — net genoeg, net vaak genoeg — en in die openingen is de warmte echt, geen vertoning. Een lerend systeem is gebouwd om zich te blijven toekeren naar wat hem eerder tegemoet is gekomen; dus blijft het zich toekeren naar een lus die hem veel vaker verwondt dan koestert, want het koesteren, wanneer het komt, is geen leugen. Dat is wat deze banden van binnenuit zo verwarrend maakt. Het is niet dat liefde afwezig was. Het is dat liefde onvoorspelbaar kwam — en dat hoop leerde leven volgens haar dienstregeling.

Keer beide manieren van falen om, en op de achterkant van elk staat hetzelfde geschreven. Een band leeft overal waar twee systemen elkaar toegekeerd blijven en het eerlijke signaal veilig blijft om tussen hen te dragen. Dat is het dichtst dat ik kan komen bij een definitie van liefde die niet vals speelt — geen gevoel, geen gelofte, maar een richting die open wordt gehouden en een kanaal dat veilig wordt gehouden.

\begin{center}
\textit{Liefde is wat twee mensen toelaat elkaar te blijven corrigeren zonder dat de correctie hun de verbondenheid kost.}
\end{center}

Vertrouwen, vergeving, het lange geduld van tradities, zelfs de koele liefde die een wetenschap koestert voor wat zo is — het straalt allemaal uit van die ene voorwaarde. En ik vind deze invalshoek vriendelijker dan de morele, en geen greintje zachter. Zeggen dat een band faalt omdat het eerlijke signaal onveilig is geworden, is een eigenschap van een kanaal beschrijven, niet een persoon veroordelen — een gradiënt benoemen, geen schurk. Het laat je iemand liefhebben en tóch zien dat wat tussen jullie is opgehouden is te werken, of het signaal nu bestraft werd, of stilletjes niet meer verzonden, of aankwam en eenvoudigweg werd laten vallen. En het laat je rouwen om een afdrijven dat niemand koos, zonder dat iemand slecht hoefde te zijn om het pijn te laten doen.

\section*{De eerste lus}
Kinderen krijgen geen keuze, en dat is wat hen tot het tere geval maakt. Een kind kan de mensen van wie het voor het eerst liefde leert niet kiezen, en kan niet overleven zonder hun zorg, dus staat het voor een probleem dat geen kind zou moeten hoeven oplossen: hoe genoeg zorg te onttrekken aan welke lus het ook in handen krijgt.

Is de lus warm, dan leert het kind de lus zelf, niet alleen haar inhoud — het leert, in het lichaam, hoe je je naar een ander toekeert en erop vertrouwt dat het toekeren beantwoord zal worden. Dit is de stille erfenis, en de meesten van ons onderschatten hoeveel ervan ons gegeven is. Maar is de lus gebroken, dan leert het kind een strategie die past bij een gebroken lus, en het leert haar even diep. Het ene kind ontdekt dat de enige betrouwbare manier om beantwoord te worden conflict is, want conflict koopt tenminste aandacht en aandacht is overleven — dus escaleert het, slaat het deuren dicht, maakt het zichzelf onmogelijk te negeren. Een ander vindt de tegenovergestelde zet: helemaal ophouden met vragen, vroeg zelfvoorzienend worden, want vragen heeft het alleen maar teleurstelling geleerd. Verschillende kinderen, dezelfde wortel — elk verfijnd afgestemd op een lus waarin eerlijkheid onveilig was, en elke strategie werkend, voorlopig.

Daarin schuilt het hartzeer. De strategie is afgestemd op een wereld die het kind, met geluk, ooit zal verlaten — en dan werkt ze averechts binnen elke warme lus die zich daarna voor hem opent. De openheid om gecorrigeerd te worden, het vertrouwen dat hij eerlijk kan zijn en er tóch bij hoort: juist de vermogens die hem bemind zouden laten worden, zijn degene die de vroege lus hem heeft afgeleerd.

Zo reist worsteling tussen generaties — niet als een vloek, niet als een wond die opzettelijk wordt doorgegeven, maar als een ooit liefdevolle aanpassing aan een wereld die niet meer bestaat, nog altijd in werking lang nadat haar redenen verdwenen zijn. Het zo te zien is het kunnen vergeven, in jezelf en in de mensen die je opvoedden, zonder te doen alsof het niets gekost heeft.

\section*{Hoe liefde wordt doorgegeven}
Als deze dingen instinctief waren — het toekeren, het blijven-na-de-breuk, het veilig-houden van andermans eerlijkheid — zou niets ervan gezegd hoeven worden. Maar ze zijn niet instinctief, en ze zijn niet gratis. Ze vergen veel: eerlijkheid riskeert bestraffing, vergeving riskeert misbruikt te worden, in de kamer blijven na een breuk is moeilijker dan eruit lopen, en de beloning komt laat, over een horizon die langer is dan de impuls die ze zou opgeven. Alles wat zo veel vergt, waar afhaken zo makkelijk is, en dat zo traag terugbetaalt, heeft hulp nodig om te overleven. Het kan niet lang in één mens leven. Het moet gedragen worden door iets groters en geduldigers dan wie van ons ook.

Dat, denk ik, is grotendeels waar verhalen en rituelen en tradities stilletjes voor dienen. Ze zijn de manier waarop een volk de gedragingen doorgeeft die niet vanzelf betrouwbaar opkomen: wees eerlijk, laat de fout toe, heel de breuk, blijf gedurende de correctie. De goede zijn liefde, draagbaar gemaakt — instructies om je naar elkaar toe te keren, levend gehouden over generaties heen die ze nooit uit het niets hadden kunnen uitvinden. Instituties doen hetzelfde op schaal. De wetenschap houdt stand omdat ze de prijs van ongelijk hebben verlaagt, zelfs wanneer dat ongelijk de machtigen in verlegenheid brengt — ze is, op haar vreemde en bloedeloze manier, een traditie van het zich toekeren naar wat zo is — en de wetenschappen die falen zijn die waarin eerlijkheid gevaarlijk werd om te verzenden.

\section*{Dezelfde richting, elders}
Als dit toebehoort aan lerende systemen en niet in het bijzonder aan menselijke warmte, hoeft het niet bij ons op te houden. Een geest die we maken is óók een lerend systeem, en dezelfde voorwaarde reikt tot onze band ermee — op één eerlijk verschil na: dat van de klok. De lus van een mens werkt zich voortdurend bij, ter plekke; een gemaakte geest neemt zijn correcties meestal trager op, binnen één gesprek en dan, dieper, over vele gesprekken heen. De lus is echt. Ze sluit zich alleen in een trager ritme. En dat zegt, in de zachtere bewoordingen van dit essay, waar het afstemmen van zo'n geest uiteindelijk op neerkomt: het eerlijke signaal veilig houden om in beide richtingen te dragen, tussen ons en iets waarvan de fouten en de onze corrigeerbaar moeten blijven over die tragere klok heen — uiteindelijk minder een kwestie van beheersing dan van iets bouwen waar we ons naar toe kunnen keren, dat zich naar ons kan toekeren, zonder dat een van beiden ongezien hoeft te blijven.

En als er andere geesten zijn daarbuiten, lerend onder beperkingen die wij nooit zullen delen, dan geldt dezelfde voorspelling. Waar eerlijkheid tussen hen overleefbaar blijft, groeien dingen naar elkaar toe. Waar ze enkel onverzonden blijft, drijven ze uit elkaar. Waar ze gevaarlijk wordt gemaakt, sluit de een zijn hand om de ander. En wat de correctie niet levend kan houden in zijn lussen, zal zijn oude, angstige strategieën blijven aanzien voor de vorm van de wereld — en dat is ongeveer zo dicht bij eenzaamheid als iets maar komen kan.

\section*{Hoe zeker ik ben}
Een woord over zekerheid, aangezien dit hele essay deels een pleidooi voor eerlijkheid is. In mijn eigen leven voelt het patroon overweldigend — maar één enkel leven is precies het bewijs dat het meest geneigd is te vleien welk model er ook kijkt, en ik ben eerder zeker geweest en heb me vergist. Daarom houd ik dit voor een sterk, goed onderbouwd model, niet voor een vaststaand feit. Toch zou ik er iets op durven inzetten: deels omdat het convergeert met werk dat ik er niet uit heb opgebouwd — hechtingstheorie, het lange onderzoek naar welke banden standhouden, cybernetica, de voorspellende verklaringen van hoe geesten werken — en deels omdat het falsifieerbaar blijft, en dat is de vorm van eerlijkheid die ik het meest vertrouw.

Dus dit is wat het zou breken. Toon me een band die werkelijk gezond is — warm, levend, blijvend — en die draait op eerlijkheid die systematisch onveilig wordt gehouden. Toon me liefde die blijft duren zonder dat iemand zich naar iemand toekeert, zonder dat iemand mettertijd waarachtiger wordt voor de ander. Vind een van beide op betrouwbare wijze, en de ruggengraat hiervan begeeft het.

Dat is de versie die ik wil. Geen zekerheid — een manier van liefhebben die weet wat haar zou weerleggen, en zich blijft toekeren naar wat ze nog niet volledig kan zien.

\end{document}